\subsection{A qualitative bound depending only on the inner dimension}

\begin{proposition}[P-F18: finite dimension envelope for ordinary plethysm]
\label{prop:innerdimensionbound}
For every integer $m\ge1$ there is a finite integer $C(m)\ge1$ with
the following property. Let $n\ge1$, let $\mu$ be a partition of
length at most $n$, and suppose
\[
 W=S^\mu\mathbb C^n,\qquad \dim W=m.
\]
For partitions $\lambda,\pi$ of lengths at most $m,n$, respectively,
put
\[
 f_{n,\mu,\lambda,\pi}(s)
 =\dim\operatorname{Hom}_{GL_n}
   \bigl(S^{s\pi}\mathbb C^n,S^{s\lambda}(W)\bigr).
 \tag{P-F18a}
\]
If this function has a nonzero positive-degree value, the canonical
graded multiplicity algebra has a positive Hilbert form whose
denominator weights are at most $C(m)$; its canonical modular index
$\chi$, defined by a lexicographically minimal h.s.o.p., also satisfies
\[
 \chi\le C(m),\qquad
 s(f_{n,\mu,\lambda,\pi})\le
 p(f_{n,\mu,\lambda,\pi})\le C(m).
 \tag{P-F18b}
\]
For a ray zero in every positive degree, assign indices and $\chi$
equal to zero. Its literal value at degree zero is one and is treated
separately from its zero quasi-polynomial constituents.

No polynomial growth, effective size estimate, or uniform algorithm
for $C(m)$ is asserted. More precisely, define the optimal envelopes
over the ordinary plethysm family (P-F18a) of inner dimension $m$ by
\[
 \Xi(m)=\max\bigl(\{1\}\cup\{\chi\}\bigr),\qquad
 P(m)=\max\bigl(\{0\}\cup\{p(f_{n,\mu,\lambda,\pi})\}\bigr).
 \tag{P-F18c}
\]
These maxima are finite. The restriction of GCT VI,
Hypothesis 1.6.7 (including its canonical modular-index clause), to
ordinary $GL_n$ plethysm is equivalent to polynomial growth of
$\Xi(m)$. The corresponding restriction of Hypothesis 1.6.6 is
equivalent to polynomial growth of $P(m)$.
Here the canonical modular index uses the parameter ring $T$
in the detailed definition of GCT VI, Section 4.1.1, printed p.~68,
and the multiplicity-ring formula (4.5), printed p.~69.
The introductory paragraph before Hypothesis 1.6.7 instead writes
$S$; we follow the detailed $T=S^H$ definition, whose Hilbert
function is the multiplicity, consistently with Theorem 1.6.1(c).
\end{proposition}

\begin{proof}
Assume first $m>1$, pad $\mu$ with zeros to length $n$, and write
\[
 k=\mu_n,\qquad \nu=\mu-k(1^n),\qquad
 \delta_i=\mu_i-\mu_{i+1}\quad(1\le i<n).
\]
Then $\nu_n=0$, the integers $\delta_i$ are nonnegative, and
\[
 S^\mu\mathbb C^n\simeq
 (\det)^k\otimes S^\nu\mathbb C^n.
 \tag{P-F18d}
\]
In particular, the dimension and the restriction to $SL_n$ are
unchanged by replacing $\mu$ with $\nu$.

We give explicit bounds on the possible normalized highest weights.
The type-$A$ Weyl dimension formula is
\[
 m=\prod_{1\le i<j\le n}
       \frac{\nu_i-\nu_j+j-i}{j-i}
   =\prod_{1\le i<j\le n}
       \left(1+\frac{\delta_i+\cdots+\delta_{j-1}}{j-i}\right).
 \tag{P-F18e}
\]
The primary locator is Borcherds--Haiman--Reshetikhin--Serganova--
Johnson-Freyd, \emph{Lectures in Lie Groups and Quantum Groups},
Part I (2016), Theorem 6.1.1.30, printed pp.~102--103;
see also its type-$A$ Example 6.1.1.31, printed p.~103.
(P-F18e) is its $GL_n$ specialization, whose central character
does not affect dimension. Every factor in (P-F18e) is at least
one, and its factor for $(i,j)=(i,i+1)$ is $\delta_i+1$.
Consequently
\[
 0\le\delta_i\le m-1\qquad(1\le i<n).
 \tag{P-F18f}
\]
Moreover $n\ge2$ and at least one $\delta_i$ is nonzero: if all
were zero, (P-F18e) would give $m=1$.
Choose an index $i$ with $\delta_i\ge1$. Every factor in (P-F18e)
is monotone in every $\delta_j$, so comparison with the weight
having only $\delta_i=1$ gives
\[
 m\ge\prod_{a=1}^{i}\prod_{b=i+1}^{n}
                  \frac{b-a+1}{b-a}
   =\prod_{a=1}^{i}\frac{n-a+1}{i-a+1}
   =\binom ni\ge n.
 \tag{P-F18g}
\]
For the last inequality, the binomial coefficients increase from
$\binom n1=n$ up to the middle because
$\binom n{j+1}/\binom nj=(n-j)/(j+1)\ge1$ there, and symmetry
handles the remaining indices $1\le i\le n-1$.

It follows that, for a fixed $m>1$, the pairs $(n,\nu)$ that can
occur belong to the finite set specified by
\[
 2\le n\le m,\qquad \nu_n=0,\qquad
 \nu_i-\nu_{i+1}\in\{0,1,\ldots,m-1\},\qquad
 \dim S^\nu\mathbb C^n=m.
 \tag{P-F18h}
\]
Indeed $\nu_i=\sum_{j=i}^{n-1}\delta_j$ is uniquely determined
by the bounded tuple of differences. There are at most
$\sum_{n=2}^{m}m^{n-1}$ tuples before imposing the dimension
condition. This is a finiteness observation, not a polynomial bound
on the representation-theoretic constants below.

Each pair in (P-F18h) defines one algebraic homomorphism
\[
 \overline\rho_{n,\nu}:SL_n\longrightarrow SL_m,
 \tag{P-F18i}
\]
after a choice of basis. The image lies in $SL_m$ because the
determinant of a representation of a connected semisimple group is
a trivial character. The map is not required to be injective;
it can have a central kernel. Changing the basis conjugates the
homomorphism and gives isomorphic canonical multiplicity algebras.
By (P-F18d), every $\mu$ of inner dimension $m$ restricts to one
of these finitely many $SL_n$ representations.

We spell out why determinant twists and the central characters do
not reintroduce dependence on their unbounded sizes. A scalar
$zI_n$ acts on the two representations in (P-F18a) by the
characters $z^{s|\pi|}$ and $z^{s|\mu||\lambda|}$.
If
\[
 |\pi|\ne |\mu||\lambda|,
 \tag{P-F18j}
\]
then all positive-degree multiplicities are zero. If equality
holds, the central actions agree, and $GL_n$-equivariance is
equivalent to $SL_n$-equivariance. Indeed every complex invertible
matrix is a product of a scalar matrix and a matrix of determinant
one. Write $\bar\lambda,\bar\pi$ for the corresponding highest
weights of $SL_m,SL_n$. On a balanced ray (P-F18a) is therefore
\[
 \dim\operatorname{Hom}_{SL_n}
  \left(V_{s\bar\pi},
       V_{s\bar\lambda}|_{\overline\rho_{n,\nu}}\right).
 \tag{P-F18k}
\]
P-F17 proves this identification as an isomorphism of the actual
graded Cartan multiplicity algebras: matrix-coefficient restriction
respects products, and on the balanced diagonal the scalar action
is trivial, so the $GL_n$ and $SL_n$ invariant subalgebras coincide.
The external grading is retained even when a highest weight is a
determinant character. Thus the conclusion applies to the canonical
modular index as well as to the numerical Hilbert function.

For each of the finitely many homomorphisms (P-F18i), choose the
finite bounds $L_{n,\nu}$ and $C_{n,\nu}$ provided by P-F17.
Take
\[
 C(m)=\max\bigl(\{1\}\cup
       \{L_{n,\nu},C_{n,\nu}:(n,\nu)\text{ satisfies (P-F18h)}\}\bigr).
 \tag{P-F18l}
\]
P-F17 now gives (P-F18b) for every balanced ray with a nonzero
positive-degree value, uniformly over $\lambda,\pi$ and every
determinant twist $k$. It also handles the trivial joint $SL$
weight, whose canonical section algebra is $\mathbb C[t]$.
The remaining rays have the stated zero convention.

If $m=1$, (P-F18e), or the single-coordinate case $n=1$, shows
that $\mu=(k^n)$ and $W=(\det)^k$. Since $\lambda$ has at most
one part,
\[
 S^{s\lambda}(W)=(\det)^{sk|\lambda|}.
\]
Thus every positive-degree multiplicity is one if
$\pi=(k|\lambda|)^n$, and zero otherwise. In the first case the
canonical algebra is $\mathbb C[t]$, its modular index is one,
and its positivity and saturation indices are zero. The second
case has the stated zero convention. Taking $C(1)=1$ completes
the finiteness proof.

The maxima in (P-F18c) exist because their entries are nonnegative
integers bounded by $C(m)$; the sets are nonempty by the added
values one and zero. These envelopes depend on the optimal
canonical modular and positivity indices, not on an arbitrary
choice of the coarse bounds in (P-F18l).
For $m>1$, (P-F18g) and the admissible partition lengths give
\[
 m^2\le
 h=\dim GL(W)+\operatorname{ht}(\lambda)+\operatorname{ht}(\pi)
 \le m^2+m+n\le m^2+2m.
 \tag{P-F18m}
\]
These are exactly the parameters in GCT VI, Hypotheses 1.6.6--1.6.7.
If the canonical modular index is bounded by a fixed polynomial in
$h$, then (P-F18m) bounds $\Xi(m)$ by a fixed polynomial in $m$.
Conversely, a polynomial bound on $\Xi(m)$ bounds every canonical
modular index polynomially in $h$, since $m^2\le h$.
The $m=1$ case satisfies both bounds directly. This proves the
claimed equivalence for the full canonical clause of ordinary
plethysm PH3; the associated positive form is supplied by its
h.s.o.p. The identical argument with $p$ proves the PH2 equivalence
for $P(m)$.
\end{proof}

\paragraph{Quantitative scope.}
This closes a qualitative uniformity question for ordinary
$GL_n$ plethysm: even allowing $n$, the determinant twist of $\mu$,
and the outer and target partitions to vary, fixed inner dimension
permits finite bounds. The canonical clause of PH3 would require
polynomial growth of the optimal envelope $\Xi(m)$. The next
result, P-F19, shows that this requirement fails on an explicit
subsequence, while the finite bound proved here remains valid.
The weaker existence of some bounded positive rational form and the
corresponding PH2 question for $P(m)$ are separate assertions.
P-F31 later refutes the former by an actual ordinary Hilbert-series
pole; the PH2 question remains open here. P-F12 supplies the valid
conditional implication between those two bounds.
The proof here neither estimates the finite bound nor constructs
the required parameters uniformly. Polynomial growth of one particular coarse choice
$C(m)$ would suffice, but is not asserted to be equivalent to PH3.
The unrestricted reductive-group formulation of Problem 1.1.2 is
not reduced to this ordinary $GL_n$ case here. For $m>1$, the parameter
$h\asymp m^2$ is also not the binary encoding length of the inner
partition, so no compressed-input polytope construction follows.
No historical priority is claimed.
